Memory is the ability to use past experience to influence the present. It relies on neuronal imprints of memories, likely in the form of long-term synaptic changes \citep{Martin2002} and neuronal activity patterns that are reactivated during memory recall \citep{Josselyn2020}. The long-term maintenance of memories over months and years requires an active process---memory consolidation---that transforms recent, malleable memories into a stable format \citep{Squire2015}. Memory consolidation occurs both at the synaptic level in the form of molecular changes that increase synaptic stability \citep{Reymann2007, Fusi2005, Clopath2012} and at the systems level in the form of a redistribution of memories across different areas of the brain \citep{Squire1995}. Both forms are thought to address the plasticity-stability dilemma, i.e., the challenge of both learning rapidly and retaining memories for long periods of time \citep{McClelland1995, Fusi2005, Roxin2013, Remme2021}. 

Systems consolidation redistributes memories across different brain regions, but the accompanying changes at the cellular level are not well understood. The act of remembering appears to rely on ``engram cells'' \citep{Josselyn2020, Han2009, Liu2012} that are consistently reactivated during memory recall over long time spans after memory acquisition. On the other hand, memories that are initially dependent on the hippocampus gradually lose this dependence, suggesting that the engram should change and gradually transform into a brain-wide ``engram complex'' of linked ensembles of engram cells \citep{Roy2022}, potentially involving the spontaneous reactivation of engrams during sleep \citep{Diekelmann2010}. 

While the dynamics of engrams as the neuronal representation of memories remain unclear \citep{Zaki2024}, the dynamics of neuronal representations of space and other task variables have recently been studied extensively. Even after a task has been successfully learned, the associated representations continue to change over the course of days to weeks \citep[e.g.,][]{Ziv2013,Driscoll2017,Rule2019,Schoonover2021}. This phenomenon, termed ``representational drift'', is experience-dependent \citep{Geva2023, Khatib2023} and therefore also relies at least partially on neuronal reactivation. Mechanistic theories of representational drift have suggested that it results from continued learning \citep{Driscoll2022a,Devalle2025} or stochastic synaptic fluctuations \citep{Qin2023, Eppler2026}. Yet, whether drift is merely a stochastic epiphenomenon that can be compensated for \citep{Rule2022, Micou2024} or whether it is a directed, purposeful process is unclear \citep{Micou2023}.

The goal of the present study is to provide a conceptual bridge between the classical, area-centered view of systems consolidation and the cellular level of neuronal engrams. To this end, we develop a phenomenological model for the dynamics of neuronal engrams during consolidation. The model is mathematically similar to a recurrent neuronal network, allowing us to reinterpret various dynamical phenomena known from neural networks in the domain of memory consolidation. It provides an engram-based view of psychological phenomena, such as selective forgetting and memory semantization, and reproduces power-law forgetting without the need for neuron- or area-specific learning rates \citep{Roxin2013}. Finally, we show that the dynamics of memory consolidation are accompanied by representational changes that share many hallmarks of representational drift, suggesting that representational drift may be a side effect of an active redistribution of memories for long-term maintenance.
